% GENERATED FILE. Edit canonical lessons, then run npm run generate:books.
% canonical-source-hash: fnv1a64:7c3b1cb0cf5a7e8c
% canonical-lessons: ES-C43-comida

\chapter{La Comida --- A Noun You Can Build}
\label{ch:noun-comida}
\hlchaptermodality{68}

\begin{chapteropening}
\textbf{By the end of this chapter:} \emph{I can build comida from comer with -ida, and know that nouns made this way are feminine.}

Complete ES-C43-comida: derive comida from comer and say what gender the -ida ending guarantees.
\end{chapteropening}

\section[la comida]{la comida — \textquotedblleft{}food,\textquotedblright{} built from a verb you have}

\label{lesson:ES-C43-comida}



\begin{quote}
Say \textquotedblleft{}to eat.\textquotedblright{} (\emph{Comer}.) Keep it in mind — this chapter's noun is made out of it, in front of you.
\end{quote}

\begin{sounds}
\begin{itemize}
\raggedright
  \item \texttt{d-soft} — the \emph{d} between vowels is soft, close to the \emph{th} in English \emph{this}: \emph{ko-MEE-thah}, not a hard \emph{d}.
\end{itemize}
\end{sounds}

\begin{cousinweb}
\begin{quote}
\textbf{la comida} — \emph{food}, and also \emph{a meal}.
\end{quote}

Take \emph{comer}, drop the ending, add \textbf{-ida}:

\begin{quote}
\textbf{com}er $\to$ \textbf{com}ida — literally \emph{the eaten thing}.
\end{quote}

That ending is not decoration. It is a working tool: it takes a verb and hands you back a noun for the \textbf{result} of that verb. You will meet it again and recognise it when you do.

This is the first Spanish noun you did not have to be given. You could have built it. That is a different kind of possession from memorising a list, and it is the reason the ending is worth more here than the word.

\emph{Comer} itself comes from Latin \textbf{comedere} — \emph{ed-} is \textquotedblleft{}eat,\textquotedblright{} the same ancient root as English \textbf{eat}, with \emph{com-} (\textquotedblleft{}thoroughly\textquotedblright{}) in front. So \emph{comida} is, at the bottom, \textquotedblleft{}the thoroughly-eaten thing.\textquotedblright{}
\end{cousinweb}

\begin{grammarlens}[title={Grammar Lens: -ida nouns are feminine}]
Nouns built this way end in \emph{-a} and are \textbf{feminine}, every time. So the ending gives you two things at once: the word, and its gender.

\begin{quote}
\emph{la comida} — the food. \emph{Hago la comida.} — I make the meal.
\end{quote}

Set your three nouns side by side and the pattern is doing all the work:

\noindent
\begin{tabularx}{\linewidth}{@{}>{\raggedright\arraybackslash}X>{\raggedright\arraybackslash}X>{\raggedright\arraybackslash}X@{}}
\toprule
\textbf{noun} & \textbf{gender} & \textbf{article} \\
\midrule
\emph{libro} & masculine & \textbf{el} \\
\emph{casa} & feminine & \textbf{la} \\
\emph{comida} & feminine & \textbf{la} \\
\bottomrule
\end{tabularx}
\end{grammarlens}

\subsection*{Your turn}
Say these aloud:

\begin{itemize}
\raggedright
  \item \textquotedblleft{}comer\textquotedblright{} then \textquotedblleft{}comida\textquotedblright{} — hear the verb inside the noun
  \item \textquotedblleft{}Hago la comida\textquotedblright{}
\end{itemize}

\begin{tcolorbox}[breakable,colback=teal!4,colframe=teal!35!black,title={Before you move on}]
What verb is hiding inside \emph{comida}? (\emph{\textbf{Comer}}.) What does the \emph{-ida} ending give you besides the word? (Its \textbf{gender} — always feminine.) Next: what happens to these three nouns, and to their articles, when there is more than one of something.
\end{tcolorbox}
